\documentclass[10pt,aspectratio=169]{beamer}
\input{../../_en_preambulo_beamer}
% Comandos y macros estadísticas compartidas para las presentaciones Beamer en inglés (Metropolis)

% Conjuntos numéricos básicos
\newcommand{\R}{\mathbb{R}}
\newcommand{\N}{\mathbb{N}}
\newcommand{\Q}{\mathbb{Q}}
\newcommand{\Z}{\mathbb{Z}}
\newcommand{\set}[1]{\left\{ #1 \right\}}
\newcommand{\abs}[1]{\left|#1\right|}
\newcommand{\norm}[1]{\left\| #1 \right\|}

% Comandos estadísticos y probabilísticos del curso
\newcommand{\Var}{\operatorname{Var}}
\newcommand{\cov}{\operatorname{Cov}}
\newcommand{\corr}{\rho}
\newcommand{\comb}[2]{\begin{pmatrix} #1 \\ #2 \end{pmatrix}}
\newcommand{\s}{\sigma}
\newcommand{\particion}{\mathfrak{P}}
\newcommand{\card}[1]{\#\left( #1 \right)}
\newcommand{\seq}[3]{\set{#1}_{#2}^{#3}}
\newcommand{\E}{\mathbb{E}}
\newcommand{\Prob}{\mathbb{P}}

% Entornos pedagógicos y cajas en inglés académico natural (soportando tanto nombres en inglés como en español)
\newenvironment{discussion}{%
  \begin{alertblock}{Discussion Question}%
}{%
  \end{alertblock}%
}
\newenvironment{discusion}{%
  \begin{alertblock}{Discussion Question}%
}{%
  \end{alertblock}%
}

\newenvironment{reflection}{%
  \begin{alertblock}{Classroom Reflection}%
}{%
  \end{alertblock}%
}
\newenvironment{reflexion}{%
  \begin{alertblock}{Classroom Reflection}%
}{%
  \end{alertblock}%
}

\newenvironment{paradox}{%
  \begin{alertblock}{Exploratory Paradox}%
}{%
  \end{alertblock}%
}
\newenvironment{paradoja}{%
  \begin{alertblock}{Exploratory Paradox}%
}{%
  \end{alertblock}%
}

\newenvironment{motivation}{%
  \begin{exampleblock}{Motivation: Data Science in Practice}%
}{%
  \end{exampleblock}%
}
\newenvironment{motivacion}{%
  \begin{exampleblock}{Motivation: Data Science in Practice}%
}{%
  \end{exampleblock}%
}

\newenvironment{quickchallenge}{%
  \begin{block}{Quick Team Challenge}%
}{%
  \end{block}%
}
\newenvironment{retoclase}{%
  \begin{block}{Quick Team Challenge}%
}{%
  \end{block}%
}


\title[Sample Size]{Section 6.7: Sample Size Estimation}
\subtitle{Confidence, margin of error, and prior variability}
\author[J. Castillo Colmenares]{Juliho Castillo Colmenares}
\institute{Tecnológico de Monterrey}
\date{}

\begin{document}
\begin{frame}[plain]\vspace{-0.8cm}\titlepage\end{frame}

\begin{frame}{Roadmap --- Chapter 6}
  \footnotesize
  \begin{itemize}\itemsep=0.08cm
    \item \textbf{06.02} Confidence intervals
    \item \textbf{06.04--06.05} Standard errors and proportions
    \item \textbf{06.07} Sample size \emph{(today)}
    \item \textbf{07} Hypothesis tests and power
  \end{itemize}
  \begin{block}{Learning objective}
    Determine $n$ before collecting data by fixing confidence, margin of error,
    and a prior variability estimate.
  \end{block}
\end{frame}

\begin{frame}{Motivation: how much data?}
  \footnotesize
  A sample that is too small may lack power; an excessively large sample wastes
  money, time, and computational infrastructure.
  \begin{alertblock}{Guiding question}
    What size guarantees that the margin of error does not exceed $E$?
  \end{alertblock}
\end{frame}

\begin{frame}{Three prior decisions}
  \footnotesize
  Before data collection specify:
  \begin{enumerate}\itemsep=0.12cm
    \item nominal confidence $(1-\alpha)$;
    \item maximum tolerable margin $E$;
    \item prior dispersion $\sigma$ or expected proportion $p^*$.
  \end{enumerate}
\end{frame}

\begin{frame}{Sample size for a mean}
  \footnotesize
  Since $E=Z_{\alpha/2}\sigma/\sqrt n$,
  \[
    n=\left(\frac{Z_{\alpha/2}\sigma}{E}\right)^2.
  \]
  If $\sigma$ is unknown, estimate it from a pilot or the range rule
  $\sigma\approx\text{range}/4$.
\end{frame}

\begin{frame}{Sample size for a proportion}
  \footnotesize
  Since $E=Z_{\alpha/2}\sqrt{p(1-p)/n}$,
  \[
    n=\left(\frac{Z_{\alpha/2}}{E}\right)^2p^*(1-p^*).
  \]
  Always round the result upward.
\end{frame}

\begin{frame}{Conservative case and finite population}
  \footnotesize
  Without prior information, $p^*=0.5$ maximizes $p(1-p)$:
  \[
    n_{\max}=\frac{Z_{\alpha/2}^2}{4E^2}.
  \]
  If $n/N>0.05$, adjust for a finite population:
  \[
    n_a=\frac{n}{1+(n-1)/N}.
  \]
\end{frame}

\begin{frame}{Monotonicity of required size}
  \footnotesize
  Size increases with $Z_{\alpha/2}$ and variability, and decreases when a
  larger margin $E$ is tolerated.
  \begin{alertblock}{Reading}
    The formulas make the trade-off between precision and resources explicit.
  \end{alertblock}
\end{frame}

\begin{frame}{Computational bridge 1/4: confidence effect}
  \footnotesize
  A higher confidence level increases $Z_{\alpha/2}$ and therefore the required
  sample size.
  \begin{block}{Rule}
    More confidence or a smaller $E$ requires more observations.
  \end{block}
\end{frame}

\begin{frame}{Computational bridge 2/4: margin effect}
  \footnotesize
  \[
    n\propto\frac1{E^2}.
  \]
  Halving the margin approximately quadruples the sample size.
\end{frame}

\begin{frame}{Computational bridge 3/4: variability}
  \footnotesize
  \begin{columns}[T]
    \column{0.48\textwidth}
      Mean: $n\propto\sigma^2$.
    \column{0.48\textwidth}
      Proportion: $n\propto p^*(1-p^*)$.
  \end{columns}
  \vspace{0.2cm}
  Reliable prior information reduces cost.
\end{frame}

\begin{frame}{Computational bridge 4/4: rounding and FPC}
  \footnotesize
  \begin{enumerate}\itemsep=0.12cm
    \item compute the continuous formula;
    \item round upward;
    \item apply FPC if the sampling fraction exceeds $5\%$.
  \end{enumerate}
\end{frame}

\begin{frame}{Reading application 1/4 --- mean}
  \footnotesize
  A prior estimate $\sigma$ is available, and confidence and $E$ are fixed.
  \begin{block}{Statement}
    Solve the mean margin-of-error equation for $n$.
  \end{block}
\end{frame}

\begin{frame}{Reading application 1/4 --- solution}
  \footnotesize
  \[
    n=\left(\frac{Z_{\alpha/2}\sigma}{E}\right)^2,
    \qquad n_{\text{final}}=\lceil n\rceil.
  \]
  Rounding upward guarantees $E_{\text{final}}\le E$.
\end{frame}

\begin{frame}{Reading application 2/4 --- proportion with $p^*$}
  \footnotesize
  A binary-response study has an expected proportion $p^*$.
  \begin{block}{Statement}
    Use anticipated variability to calculate the required size.
  \end{block}
\end{frame}

\begin{frame}{Reading application 2/4 --- solution}
  \footnotesize
  \[
    n=\left(\frac{Z_{\alpha/2}}{E}\right)^2p^*(1-p^*).
  \]
  A prior proportion far from $0.5$ implies lower expected variance.
\end{frame}

\begin{frame}{Reading application 3/4 --- conservative case}
  \footnotesize
  If no prior information about $p$ exists, protect against the worst case.
  \begin{block}{Statement}
    Maximize $p(1-p)$ on $[0,1]$.
  \end{block}
\end{frame}

\begin{frame}{Reading application 3/4 --- solution}
  \footnotesize
  \[
    p^*=0.5,\qquad p^*(1-p^*)=0.25,
    \qquad n_{\max}=\frac{Z_{\alpha/2}^2}{4E^2}.
  \]
  This size does not underestimate possible proportion variability.
\end{frame}

\begin{frame}{Reading application 4/4 --- finite population}
  \footnotesize
  A population of size $N$ may be small relative to the computed sample.
  \begin{block}{Statement}
    Check $n/N>0.05$ and adjust the size when needed.
  \end{block}
\end{frame}

\begin{frame}{Reading application 4/4 --- solution}
  \footnotesize
  \[
    n_a=\frac{n}{1+(n-1)/N}.
  \]
  The adjustment reduces observations without sacrificing the stipulated margin
  of error.
\end{frame}

\begin{frame}{Synthesis of the section}
  \footnotesize
  \begin{itemize}\itemsep=0.12cm
    \item Confidence, margin, and variability determine $n$.
    \item The margin enters squared in the denominator.
    \item $p=0.5$ is the conservative proportion case.
    \item FPC corrects finite populations.
  \end{itemize}
\end{frame}

\begin{frame}{Curricular transition}
  \footnotesize
  \begin{block}{Established}
    Sample size is decided before observing data, with explicit assumptions and
    safe rounding.
  \end{block}
  \begin{alertblock}{Next unit}
    Hypothesis tests use these sizes to control significance, power, and
    precision.
  \end{alertblock}
\end{frame}
\end{document}
